\section{L}

% lachen (to laugh)
\begin{verbentry}{
  style = {default},
  toc = {lachen (to laugh)},
  name = {lachen},
  english = {to laugh},
  type = {regular},
  auxiliary = {haben}
}
 

\vspace{5pt}
\hrule
\vspace{5pt}

\begin{tabular}{rl}
\multicolumn{2}{l}{\textbf{PRÄSENS}}\\
ich & \textbf{lache}\\
du & \textbf{lachst}\\
er/sie/es & \textbf{lacht}\\
wir & \textbf{lachen}\\
ihr & \textbf{lacht}\\
sie/Sie & \textbf{lachen}\\
\end{tabular}
\hfill
\begin{tabular}{rl}
\multicolumn{2}{l}{\textbf{PRÄTERITUM}}\\
ich & \textbf{lachte}\\
du & \textbf{lachtest}\\
er/sie/es & \textbf{lachte}\\
wir & \textbf{lachten}\\
ihr & \textbf{lachtet}\\
sie/Sie & \textbf{lachten}\\
\end{tabular}
\hfill
\begin{tabular}{rl}
\multicolumn{2}{l}{\textbf{IMPERATIV}}\\
& \textbf{lach(e) (du)!}\\
& \textbf{lachen wir!}\\
& \textbf{lacht (ihr)!}\\
& \textbf{lachen Sie!}\\
\end{tabular}

\vspace{10pt}

\begin{tabular}{rl}
\multicolumn{2}{l}{\textbf{PERFEKT}}\\
& haben + \textbf{gelacht}\\
\end{tabular}
\hfill
\begin{tabular}{rl}
\multicolumn{2}{l}{\textbf{FUTURE I}}\\
& werden + \textbf{lachen}\\
\end{tabular}
\hfill
\begin{tabular}{rl}
\multicolumn{2}{l}{\textbf{PLUSQUAMPERFEKT}}\\
& hatten + \textbf{gelacht}\\
\end{tabular}

\vspace{10pt}

\begin{tabular}{lll}
\multicolumn{3}{l}{\textbf{MIT PRÄPOSITIONEN}}\\
& \textbf{lachen über} + Akk & (to laugh about)\\
\end{tabular}
\hfill
\begin{tabular}{lll}
\multicolumn{3}{l}{\textbf{TRENNBARE VERBEN}}\\
& aus$|$lachen & (to laugh at / mock)\\
\end{tabular}
\vspace{-5pt}
\end{verbentry}

% lächeln (to smile)
\begin{verbentry}{
  style = {default},
  toc = {lächeln (to smile)},
  name = {lächeln},
  english = {to smile},
  type = {regular},
  auxiliary = {haben}
}


\vspace{5pt}
\hrule
\vspace{5pt}

\begin{tabular}{rl}
\multicolumn{2}{l}{\textbf{PRÄSENS}}\\
ich & \textbf{lächele}\\
du & \textbf{lächelst}\\
er/sie/es & \textbf{lächelt}\\
wir & \textbf{lächeln}\\
ihr & \textbf{lächelt}\\
sie/Sie & \textbf{lächeln}\\
\end{tabular}
\hfill
\begin{tabular}{rl}
\multicolumn{2}{l}{\textbf{PRÄTERITUM}}\\
ich & \textbf{lächelte}\\
du & \textbf{lächeltest}\\
er/sie/es & \textbf{lächelte}\\
wir & \textbf{lächelten}\\
ihr & \textbf{lächeltet}\\
sie/Sie & \textbf{lächelten}\\
\end{tabular}
\hfill
\begin{tabular}{rl}
\multicolumn{2}{l}{\textbf{IMPERATIV}}\\
& \textbf{lächle (du)!}\\
& \textbf{lächeln wir!}\\
& \textbf{lächelt (ihr)!}\\
& \textbf{lächeln Sie!}\\
\end{tabular}

\vspace{10pt}

\begin{tabular}{rl}
\multicolumn{2}{l}{\textbf{PERFEKT}}\\
& haben + \textbf{gelächelt}\\
\end{tabular}
\hfill
\begin{tabular}{rl}
\multicolumn{2}{l}{\textbf{FUTURE I}}\\
& werden + \textbf{lächeln}\\
\end{tabular}
\hfill
\begin{tabular}{rl}
\multicolumn{2}{l}{\textbf{PLUSQUAMPERFEKT}}\\
& hatten + \textbf{gelächelt}\\
\end{tabular}

\vspace{10pt}

\begin{tabular}{lll}
\multicolumn{3}{l}{\textbf{MIT PRÄPOSITIONEN}}\\
& \textbf{lächeln über} + Akk & (to smile about)\\
& \textbf{lächeln zu} + Dat & (to smile at / toward)\\
\end{tabular}
\hfill
\begin{tabular}{lll}
\multicolumn{3}{l}{\textbf{TRENNBARE VERBEN}}\\
& an$|$lächeln & (to smile at)\\
\vspace{10pt}
\end{tabular}

\vspace{-5pt}
\end{verbentry}

% laden (to load / invite)
\begin{verbentry}{
  style = {acc},
  toc = {laden (to load / invite)},
  name = {laden},
  english = {to load / invite},
  type = {irregular},
  auxiliary = {haben}
}
 
    
     \vspace{5pt}

    \hrule
    
    \vspace{5pt}

  \begin{tabular}{rl}
     \multicolumn{2}{l}{\textbf{PRÄSENS} }\\
    ich & \textbf{lade}\\
    du & \textbf{lädst}\\
    er/sie/es & \textbf{lädt}\\
    wir & \textbf{laden}\\
    ihr & \textbf{ladet}\\
    sie/Sie & \textbf{laden}\\
  \end{tabular}
  \hfill
  \begin{tabular}{rl}
    \multicolumn{2}{l}{\textbf{PRÄTERITUM} }\\
    ich & \textbf{lud}\\
    du & \textbf{ludst}\\
    er/sie/es & \textbf{lud}\\
    wir & \textbf{luden}\\
    ihr & \textbf{ludet}\\
    sie/Sie & \textbf{luden}\\
  \end{tabular}
  \hfill
  \begin{tabular}{rl}
     \multicolumn{2}{l}{\textbf{IMPERATIV} }\\
   & \textbf{lad(e) (du)!}\\
   & \textbf{laden wir!}\\
   & \textbf{ladet (ihr)!}\\
   & \textbf{laden Sie!}\\
   & \vspace{10pt}
  \end{tabular}

 \vspace{10pt}

 \begin{tabular}{rl}
     \multicolumn{2}{l}{\textbf{PERFEKT}}\\
   & haben + \textbf{geladen}\\
  \end{tabular}
\hfill
\begin{tabular}{rl}
    \multicolumn{2}{l}{\textbf{FUTURE I}}\\
   & werden + \textbf{laden}\\
  \end{tabular}
\hfill
\begin{tabular}{rl}
    \multicolumn{2}{l}{\textbf{PLUSQUAMPERFEKT}}\\
   & hatten + \textbf{geladen}\\
  \end{tabular}

 \vspace{10pt}

\begin{tabular}{lll}
      \multicolumn{3}{l}{\textbf{MIT PRÄPOSITIONEN}}\\
& \textbf{laden auf} + Akk & (to load onto)\\
& \textbf{laden in} + Akk & (to load into)\\
& \textbf{laden zu} + Dat & (to invite to)\\
\end{tabular}
  \hfill
\begin{tabular}{lll}
      \multicolumn{3}{l}{\textbf{TRENNBARE VERBEN}}\\
 & \textbf{ein$|$laden} & (to invite)\\
 & \textbf{auf$|$laden} & (to load up)\\
\vspace{10pt}
  \end{tabular}
  
  \vspace{-5pt}
\end{verbentry}

% landen (to land)
\begin{verbentry}{
  style = {default},
  toc = {landen (to land)},
  name = {landen},
  english = {to land},
  type = {regular},
  auxiliary = {sein}
}
 
    
     \vspace{5pt}

    \hrule
    
    \vspace{5pt}

  \begin{tabular}{rl}
     \multicolumn{2}{l}{\textbf{PRÄSENS} }\\
    ich & \textbf{lande}\\
    du & \textbf{landest}\\
    er/sie/es & \textbf{landet}\\
    wir & \textbf{landen}\\
    ihr & \textbf{landet}\\
    sie/Sie & \textbf{landen}\\
  \end{tabular}
  \hfill
     \begin{tabular}{rl}
    \multicolumn{2}{l}{\textbf{PRÄTERITUM} }\\
    ich & \textbf{landete}\\
    du & \textbf{landetest}\\
    er/sie/es & \textbf{landete}\\
    wir & \textbf{landeten}\\
    ihr & \textbf{landetet}\\
    sie/Sie & \textbf{landeten}\\
  \end{tabular}
  \hfill
  \begin{tabular}{rl}
     \multicolumn{2}{l}{\textbf{IMPERATIV} }\\
   & \textbf{land(e) (du)!}\\
   & \textbf{landen wir!}\\
   & \textbf{landet (ihr)!}\\
   & \textbf{landen Sie!}\\
   & \vspace{10pt}
  \end{tabular}

    \vspace{10pt}

 \begin{tabular}{rl}
     \multicolumn{2}{l}{\textbf{PERFEKT}}\\
   & sein + \textbf{gelandet}\\
  \end{tabular}
\hfill
   \begin{tabular}{rl}
   
    \multicolumn{2}{l}{\textbf{FUTURE I} }\\
   & werden + \textbf{landen}\\
 
  \end{tabular}
\hfill
   \begin{tabular}{rl}
      \multicolumn{2}{l}{\textbf{PLUSQUAMPERFEKT}}\\
   & waren + \textbf{gelandet}\\
  \end{tabular}

   \vspace{10pt}

   \begin{tabular}{lll}
      \multicolumn{3}{l}{\textbf{MIT PRÄPOSITIONEN}}\\
& \textbf{landen auf} + Dat & (to land on – surface / area)\\
& \textbf{landen in} + Dat & (to land in – country / city)\\
& \textbf{landen bei} + Dat & (to land near / at)\\
\vspace{10pt}
   \end{tabular}
  \hfill
   \begin{tabular}{lll}
      \multicolumn{3}{l}{\textbf{TRENNBARE VERBEN}}\\
 & \textbf{an$|$landen} & (to arrive / land at a place)\\
 & \textbf{auf$|$landen} & (to land on – surface)\\
 & \textbf{ein$|$landen} & (to land in / enter)\\
 & \textbf{zurück$|$landen} & (to land back / return)\\
  \end{tabular}

\vspace{-5pt}
\end{verbentry}

% laufen (to run)
\begin{verbentry}{
  style = {default},
  toc = {laufen (to run)},
  name = {laufen},
  english = {to run},
  type = {irregular},
  auxiliary = {sein}
}
 
    
     \vspace{5pt}

    \hrule
    
    \vspace{5pt}

  \begin{tabular}{rl}
     \multicolumn{2}{l}{\textbf{PRÄSENS} }\\
    ich & \textbf{laufe}\\
    du & \textbf{läufst}\\
    er/sie/es & \textbf{läuft}\\
    wir & \textbf{laufen}\\
    ihr & \textbf{lauft}\\
    sie/Sie & \textbf{laufen}\\
  \end{tabular}
  \hfill
     \begin{tabular}{rl}
    \multicolumn{2}{l}{\textbf{PRÄTERITUM} }\\
    ich & \textbf{lief}\\
    du & \textbf{liefst}\\
    er/sie/es & \textbf{lief}\\
    wir & \textbf{liefen}\\
    ihr & \textbf{lieft}\\
    sie/Sie & \textbf{liefen}\\
  \end{tabular}
  \hfill
  \begin{tabular}{rl}
     \multicolumn{2}{l}{\textbf{IMPERATIV} }\\
   & \textbf{lauf(e) (du)!}\\
   & \textbf{laufen wir!}\\
   & \textbf{lauft (ihr)!}\\
   & \textbf{laufen Sie!}\\
   & \vspace{10pt}
  \end{tabular}

    \vspace{10pt}

 \begin{tabular}{rl}
     \multicolumn{2}{l}{\textbf{PERFEKT}}\\
  &  sein + \textbf{gelaufen}\\
  \end{tabular}
\hfill
   \begin{tabular}{rl}
   
    \multicolumn{2}{l}{\textbf{FUTURE I} }\\
  &  werden + \textbf{laufen}\\
 
  \end{tabular}
\hfill
   \begin{tabular}{rl}
      \multicolumn{2}{l}{\textbf{PLUSQUAMPERFEKT}}\\
  &  waren + \textbf{gelaufen}\\
  \end{tabular}
  
  \vspace{10pt}



\begin{tabular}{lll}
\multicolumn{3}{l}{\textbf{MIT PRÄPOSITIONEN}}\\
& \textbf{---} & \\
\end{tabular}
\hfill
\begin{tabular}{lll}
\multicolumn{3}{l}{\textbf{TRENNBARE VERBEN}}\\
& \textbf{---} & \\
\end{tabular}
  
\vspace{-5pt}
\end{verbentry}

% leben (to live)
\begin{verbentry}{
  style = {default},
  toc = {leben (to live)},
  name = {leben},
  english = {to live},
  type = {regular},
  auxiliary = {haben}
}
 
    
     \vspace{5pt}

    \hrule
    
    \vspace{5pt}

  \begin{tabular}{rl}
     \multicolumn{2}{l}{\textbf{PRÄSENS} }\\
    ich & \textbf{lebe}\\
    du & \textbf{lebst}\\
    er/sie/es & \textbf{lebt}\\
    wir & \textbf{leben}\\
    ihr & \textbf{lebt}\\
    sie/Sie & \textbf{leben}\\
  \end{tabular}
  \hfill
     \begin{tabular}{rl}
    \multicolumn{2}{l}{\textbf{PRÄTERITUM} }\\
    ich & \textbf{lebte}\\
    du & \textbf{lebtest}\\
    er/sie/es & \textbf{lebte}\\
    wir & \textbf{lebten}\\
    ihr & \textbf{lebtet}\\
    sie/Sie & \textbf{lebten}\\
  \end{tabular}
  \hfill
  \begin{tabular}{rl}
     \multicolumn{2}{l}{\textbf{IMPERATIV} }\\
   & \textbf{leb(e) (du)!}\\
   & \textbf{leben wir!}\\
   & \textbf{lebt (ihr)!}\\
   & \textbf{leben Sie!}\\
   & \vspace{10pt}
  \end{tabular}

    \vspace{10pt}

 \begin{tabular}{rl}
     \multicolumn{2}{l}{\textbf{PERFEKT}}\\
   & haben + \textbf{gelebt}\\
  \end{tabular}
\hfill
   \begin{tabular}{rl}
   
    \multicolumn{2}{l}{\textbf{FUTURE I} }\\
   & werden + \textbf{leben}\\
 
  \end{tabular}
\hfill
   \begin{tabular}{rl}
      \multicolumn{2}{l}{\textbf{PLUSQUAMPERFEKT}}\\
  &  hatten + \textbf{gelebt}\\
  \end{tabular}
  
  \vspace{10pt}

\begin{tabular}{lll}
\multicolumn{3}{l}{\textbf{MIT PRÄPOSITIONEN}}\\
& \textbf{---} & \\
\end{tabular}
\hfill
\begin{tabular}{lll}
\multicolumn{3}{l}{\textbf{TRENNBARE VERBEN}}\\
& \textbf{---} & \\
\end{tabular}

  \vspace{-5pt}
\end{verbentry}

% legen (to lay / put)
\begin{verbentry}{
  style = {default},
  toc = {legen (to lay / put)},
  name = {legen},
  english = {to lay / put},
  type = {regular},
  auxiliary = {haben}
}


\vspace{5pt}
\hrule
\vspace{5pt}

\begin{tabular}{rl}
\multicolumn{2}{l}{\textbf{PRÄSENS}}\\
ich & \textbf{lege}\\
du & \textbf{legst}\\
er/sie/es & \textbf{legt}\\
wir & \textbf{legen}\\
ihr & \textbf{legt}\\
sie/Sie & \textbf{legen}\\
\end{tabular}
\hfill
\begin{tabular}{rl}
\multicolumn{2}{l}{\textbf{PRÄTERITUM}}\\
ich & \textbf{legte}\\
du & \textbf{legtest}\\
er/sie/es & \textbf{legte}\\
wir & \textbf{legten}\\
ihr & \textbf{legtet}\\
sie/Sie & \textbf{legten}\\
\end{tabular}
\hfill
\begin{tabular}{rl}
\multicolumn{2}{l}{\textbf{IMPERATIV}}\\
& \textbf{lege (du)!}\\
& \textbf{legen wir!}\\
& \textbf{legt (ihr)!}\\
& \textbf{legen Sie!}\\
\end{tabular}

\vspace{10pt}

\begin{tabular}{rl}
\multicolumn{2}{l}{\textbf{PERFEKT}}\\
& haben + \textbf{gelegt}\\
\end{tabular}
\hfill
\begin{tabular}{rl}
\multicolumn{2}{l}{\textbf{FUTURE I}}\\
& werden + \textbf{legen}\\
\end{tabular}
\hfill
\begin{tabular}{rl}
\multicolumn{2}{l}{\textbf{PLUSQUAMPERFEKT}}\\
& hatten + \textbf{gelegt}\\
\end{tabular}

\vspace{10pt}

\begin{tabular}{lll}
\multicolumn{3}{l}{\textbf{MIT PRÄPOSITIONEN}}\\
& \textbf{legen auf} + Akk & (to lay on)\\
& \textbf{legen in} + Akk & (to put into)\\
& \textbf{legen unter} + Akk & (to lay under)\\
\end{tabular}
\hfill
\begin{tabular}{lll}
\multicolumn{3}{l}{\textbf{TRENNBARE VERBEN}}\\
& \textbf{hin$|$legen} & (to lay down)\\
& \textbf{auf$|$legen} & (to put on / overlay)\\
\end{tabular}
\vspace{-5pt}
\end{verbentry}

% lehnen (to lean)
\begin{verbentry}{
  style = {default},
  toc = {lehnen (to lean)},
  name = {lehnen},
  english = {to lean},
  type = {regular},
  auxiliary = {haben / sein reflexive}
}


\vspace{5pt}
\hrule
\vspace{5pt}

\begin{tabular}{rl}
\multicolumn{2}{l}{\textbf{PRÄSENS}}\\
ich & \textbf{lehne}\\
du & \textbf{lehnst}\\
er/sie/es & \textbf{lehnt}\\
wir & \textbf{lehnen}\\
ihr & \textbf{lehnt}\\
sie/Sie & \textbf{lehnen}\\
\end{tabular}
\hfill
\begin{tabular}{rl}
\multicolumn{2}{l}{\textbf{PRÄTERITUM}}\\
ich & \textbf{lehnte}\\
du & \textbf{lehntest}\\
er/sie/es & \textbf{lehnte}\\
wir & \textbf{lehnten}\\
ihr & \textbf{lehntet}\\
sie/Sie & \textbf{lehnten}\\
\end{tabular}
\hfill
\begin{tabular}{rl}
\multicolumn{2}{l}{\textbf{IMPERATIV}}\\
& \textbf{lehne (du)!}\\
& \textbf{lehnen wir!}\\
& \textbf{lehnt (ihr)!}\\
& \textbf{lehnen Sie!}\\
\end{tabular}

\vspace{10pt}

\begin{tabular}{rl}
\multicolumn{2}{l}{\textbf{PERFEKT}}\\
& haben + \textbf{gelehnt}\\
\end{tabular}
\hfill
\begin{tabular}{rl}
\multicolumn{2}{l}{\textbf{FUTURE I}}\\
& werden + \textbf{lehnen}\\
\end{tabular}
\hfill
\begin{tabular}{rl}
\multicolumn{2}{l}{\textbf{PLUSQUAMPERFEKT}}\\
& hatten + \textbf{gelehnt}\\
\end{tabular}

\vspace{10pt}

\begin{tabular}{lll}
\multicolumn{3}{l}{\textbf{MIT PRÄPOSITIONEN}}\\
& \textbf{lehnen an} + Dat & (to lean against)\\
& \textbf{sich lehnen gegen} + Akk & (to lean against)\\
\end{tabular}
\hfill
\begin{tabular}{lll}
\multicolumn{3}{l}{\textbf{TRENNBARE VERBEN}}\\
& \textbf{---} & \\
\end{tabular}
\vspace{-5pt}
\end{verbentry}

% leihen (to borrow)
\begin{verbentry}{
  style = {acc},
  toc = {leihen (to borrow)},
  name = {leihen},
  english = {to borrow},
  type = {irregular, + Akkusativ},
  auxiliary = {haben}
}
 
    
     \vspace{5pt}

    \hrule
    
    \vspace{5pt}

  \begin{tabular}{rl}
     \multicolumn{2}{l}{\textbf{PRÄSENS} }\\
    ich & \textbf{leihe}\\
    du & \textbf{leihst}\\
    er/sie/es & \textbf{leiht}\\
    wir & \textbf{leihen}\\
    ihr & \textbf{leiht}\\
    sie/Sie & \textbf{leihen}\\
  \end{tabular}
  \hfill
     \begin{tabular}{rl}
    \multicolumn{2}{l}{\textbf{PRÄTERITUM} }\\
    ich & \textbf{lieh}\\
    du & \textbf{liehst}\\
    er/sie/es & \textbf{lieh}\\
    wir & \textbf{liehen}\\
    ihr & \textbf{lieht}\\
    sie/Sie & \textbf{liehen}\\
  \end{tabular}
  \hfill
  \begin{tabular}{rl}
     \multicolumn{2}{l}{\textbf{IMPERATIV} }\\
   & \textbf{leih(e) (du)!}\\
   & \textbf{leihen wir!}\\
   & \textbf{leiht (ihr)!}\\
   & \textbf{leihen Sie!}\\
   & \vspace{10pt}
  \end{tabular}

    \vspace{10pt}

 \begin{tabular}{rl}
     \multicolumn{2}{l}{\textbf{PERFEKT}}\\
   & sein + \textbf{geliehen}\\
  \end{tabular}
\hfill
   \begin{tabular}{rl}
   
    \multicolumn{2}{l}{\textbf{FUTURE I} }\\
   & werden + \textbf{leihen}\\
 
  \end{tabular}
\hfill
   \begin{tabular}{rl}
      \multicolumn{2}{l}{\textbf{PLUSQUAMPERFEKT}}\\
   & waren + \textbf{geliehen}\\
  \end{tabular}
  


 \vspace{10pt}

   \begin{tabular}{lll}
      \multicolumn{3}{l}{\textbf{MIT PRÄPOSITIONEN}}\\
& \textbf{leihen von} + Dat & (to borrow from)\\
& \textbf{leihen an} + Dat & (to lend to)\\
  \end{tabular}
  \hfill
   \begin{tabular}{lll}
      \multicolumn{3}{l}{\textbf{TRENNBARE VERBEN}}\\
 & \textbf{aus$|$leihen} & (to lend out (library / shop))\\
  \end{tabular}


  \vspace{-5pt}
\end{verbentry}

% lernen (to learn)
\begin{verbentry}{
  style = {acc},
  toc = {lernen (to learn)},
  name = {lernen},
  english = {to learn},
  type = {regular, + Akkusativ},
  auxiliary = {haben}
}
 
    
     \vspace{5pt}

    \hrule
    
    \vspace{5pt}

  \begin{tabular}{rl}
     \multicolumn{2}{l}{\textbf{PRÄSENS} }\\
    ich & \textbf{lerne}\\
    du & \textbf{lernst}\\
    er/sie/es & \textbf{lernt}\\
    wir & \textbf{lernen}\\
    ihr & \textbf{lernt}\\
    sie/Sie & \textbf{lernen}\\
  \end{tabular}
  \hfill
     \begin{tabular}{rl}
    \multicolumn{2}{l}{\textbf{PRÄTERITUM} }\\
    ich & \textbf{lernte}\\
    du & \textbf{lerntest}\\
    er/sie/es & \textbf{lernte}\\
    wir & \textbf{lernten}\\
    ihr & \textbf{lerntet}\\
    sie/Sie & \textbf{lernten}\\
  \end{tabular}
  \hfill
  \begin{tabular}{rl}
     \multicolumn{2}{l}{\textbf{IMPERATIV} }\\
   & \textbf{lern(e) (du)!}\\
   & \textbf{lernen wir!}\\
   & \textbf{lernt (ihr)!}\\
   & \textbf{lernen Sie!}\\
   & \vspace{10pt}
  \end{tabular}

    \vspace{10pt}

 \begin{tabular}{rl}
     \multicolumn{2}{l}{\textbf{PERFEKT}}\\
  &  haben + \textbf{gelernt}\\
  \end{tabular}
\hfill
   \begin{tabular}{rl}
   
    \multicolumn{2}{l}{\textbf{FUTURE I} }\\
  &  werden + \textbf{lernen}\\
 
  \end{tabular}
\hfill
   \begin{tabular}{rl}
      \multicolumn{2}{l}{\textbf{PLUSQUAMPERFEKT}}\\
   & hatten + \textbf{gelernt}\\
  \end{tabular}
  
  
\vspace{10pt}

\begin{tabular}{lll}
\multicolumn{3}{l}{\textbf{MIT PRÄPOSITIONEN}}\\
& \textbf{---} & \\
\end{tabular}
\hfill
\begin{tabular}{lll}
\multicolumn{3}{l}{\textbf{TRENNBARE VERBEN}}\\
& \textbf{---} & \\
\end{tabular}
\vspace{-5pt}
\end{verbentry}

% lesen (to read)
\begin{verbentry}{
  style = {acc},
  toc = {lesen (to read)},
  name = {lesen},
  english = {to read},
  type = {irregular, + acusative},
  auxiliary = {haben}
}
 
    
     \vspace{5pt}

    \hrule
    
    \vspace{5pt}

  \begin{tabular}{rl}
     \multicolumn{2}{l}{\textbf{PRÄSENS} }\\
    ich & \textbf{lese}\\
    du & \textbf{liest}\\
    er/sie/es & \textbf{liest}\\
    wir & \textbf{lesen}\\
    ihr & \textbf{lest}\\
    sie/Sie & \textbf{lesen}\\
  \end{tabular}
  \hfill
     \begin{tabular}{rl}
    \multicolumn{2}{l}{\textbf{PRÄTERITUM} }\\
    ich & \textbf{las}\\
    du & \textbf{lastest}\\
    er/sie/es & \textbf{las}\\
    wir & \textbf{lasen}\\
    ihr & \textbf{last}\\
    sie/Sie & \textbf{lasen}\\
  \end{tabular}
  \hfill
  \begin{tabular}{rl}
     \multicolumn{2}{l}{\textbf{IMPERATIV} }\\
   & \textbf{lies (du)!}\\
   & \textbf{lesen wir!}\\
   & \textbf{lest (ihr)!}\\
   & \textbf{lesen Sie!}\\
   & \vspace{10pt}
  \end{tabular}

    \vspace{10pt}

 \begin{tabular}{rl}
     \multicolumn{2}{l}{\textbf{PERFEKT}}\\
  &  haben + \textbf{gelesen}\\
  \end{tabular}
\hfill
   \begin{tabular}{rl}
   
    \multicolumn{2}{l}{\textbf{FUTURE I} }\\
  &  werden + \textbf{lesen}\\
 
  \end{tabular}
\hfill
   \begin{tabular}{rl}
      \multicolumn{2}{l}{\textbf{PLUSQUAMPERFEKT}}\\
   & hatten + \textbf{gelesen}\\
  \end{tabular}
  
  
\vspace{10pt}

\begin{tabular}{lll}
\multicolumn{3}{l}{\textbf{MIT PRÄPOSITIONEN}}\\
& \textbf{---} & \\
\end{tabular}
\hfill
\begin{tabular}{lll}
\multicolumn{3}{l}{\textbf{TRENNBARE VERBEN}}\\
& \textbf{---} & \\
\end{tabular}
\vspace{-5pt}
\end{verbentry}

% lieben (to love)
\begin{verbentry}{
  style = {acc},
  toc = {lieben (to love)},
  name = {lieben},
  english = {to love},
  type = {regular, + Akkusativ},
  auxiliary = {haben}
}
 
    
     \vspace{5pt}

    \hrule
    
    \vspace{5pt}

  \begin{tabular}{rl}
     \multicolumn{2}{l}{\textbf{PRÄSENS} }\\
    ich & \textbf{liebe}\\
    du & \textbf{liebst}\\
    er/sie/es & \textbf{liebt}\\
    wir & \textbf{lieben}\\
    ihr & \textbf{liebt}\\
    sie/Sie & \textbf{lieben}\\
  \end{tabular}
  \hfill
     \begin{tabular}{rl}
    \multicolumn{2}{l}{\textbf{PRÄTERITUM} }\\
    ich & \textbf{liebte}\\
    du & \textbf{liebtest}\\
    er/sie/es & \textbf{liebte}\\
    wir & \textbf{liebten}\\
    ihr & \textbf{liebtet}\\
    sie/Sie & \textbf{liebten}\\
  \end{tabular}
  \hfill
  \begin{tabular}{rl}
     \multicolumn{2}{l}{\textbf{IMPERATIV} }\\
   & \textbf{lieb(e) (du)!}\\
   & \textbf{lieben wir!}\\
   & \textbf{liebt (ihr)!}\\
   & \textbf{lieben Sie!}\\
   & \vspace{10pt}
  \end{tabular}

    \vspace{10pt}

 \begin{tabular}{rl}
     \multicolumn{2}{l}{\textbf{PERFEKT}}\\
   & haben + \textbf{geliebt}\\
  \end{tabular}
\hfill
   \begin{tabular}{rl}
   
    \multicolumn{2}{l}{\textbf{FUTURE I} }\\
   & werden + \textbf{lieben}\\
 
  \end{tabular}
\hfill
   \begin{tabular}{rl}
      \multicolumn{2}{l}{\textbf{PLUSQUAMPERFEKT}}\\
   & hatten + \textbf{geliebt}\\
  \end{tabular}
  
  
\vspace{10pt}

\begin{tabular}{lll}
\multicolumn{3}{l}{\textbf{MIT PRÄPOSITIONEN}}\\
& \textbf{---} & \\
\end{tabular}
\hfill
\begin{tabular}{lll}
\multicolumn{3}{l}{\textbf{TRENNBARE VERBEN}}\\
& \textbf{---} & \\
\end{tabular}
\vspace{-5pt}
\end{verbentry}

% loben (to praise)
\begin{verbentry}{
  style = {default},
  toc = {loben (to praise)},
  name = {loben},
  english = {to praise},
  type = {regular},
  auxiliary = {haben}
}
 
    
     \vspace{5pt}

    \hrule
    
    \vspace{5pt}

  \begin{tabular}{rl}
     \multicolumn{2}{l}{\textbf{PRÄSENS} }\\
    ich & \textbf{lobe}\\
    du & \textbf{lobst}\\
    er/sie/es & \textbf{lobt}\\
    wir & \textbf{loben}\\
    ihr & \textbf{lobt}\\
    sie/Sie & \textbf{loben}\\
  \end{tabular}
  \hfill
     \begin{tabular}{rl}
    \multicolumn{2}{l}{\textbf{PRÄTERITUM} }\\
    ich & \textbf{lobte}\\
    du & \textbf{lobtest}\\
    er/sie/es & \textbf{lobte}\\
    wir & \textbf{lobten}\\
    ihr & \textbf{lobtet}\\
    sie/Sie & \textbf{lobten}\\
  \end{tabular}
  \hfill
  \begin{tabular}{rl}
     \multicolumn{2}{l}{\textbf{IMPERATIV} }\\
   & \textbf{lob(e) (du)!}\\
   & \textbf{loben wir!}\\
   & \textbf{lobt (ihr)!}\\
   & \textbf{loben Sie!}\\
   & \vspace{10pt}
  \end{tabular}

    \vspace{10pt}

 \begin{tabular}{rl}
     \multicolumn{2}{l}{\textbf{PERFEKT}}\\
   & haben + \textbf{gelobt}\\
  \end{tabular}
\hfill
   \begin{tabular}{rl}
    \multicolumn{2}{l}{\textbf{FUTURE I} }\\
   & werden + \textbf{loben}\\
  \end{tabular}
\hfill
   \begin{tabular}{rl}
      \multicolumn{2}{l}{\textbf{PLUSQUAMPERFEKT}}\\
   & hatten + \textbf{gelobt}\\
  \end{tabular}

   \vspace{10pt}

   \begin{tabular}{lll}
      \multicolumn{3}{l}{\textbf{MIT PRÄPOSITIONEN}}\\
& \textbf{loben für} + Akk & (to praise for)\\
& \textbf{loben wegen} + Gen & (to praise because of)\\
\end{tabular}
  \hfill
   \begin{tabular}{lll}
\multicolumn{3}{l}{\textbf{TRENNBARE VERBEN}}\\
& \textbf{---} & \\
\end{tabular}

  \vspace{-5pt}
\end{verbentry}
